\documentclass[12pt]{article}
\usepackage[utf8]{inputenc}
\usepackage[T1]{fontenc}
\usepackage{hyperref}
\usepackage{geometry}
\usepackage{booktabs}
\usepackage{tcolorbox}
\usepackage{enumitem}
\usepackage{tabularx}

\geometry{a4paper, margin=1in}

\title{RAG and Project Architecture: Full Conversation Log}
\author{Gemini Content Generator Project}
\date{\today}

\begin{document}

\maketitle

\section{Topic 1: Current RAG Implementation}
\textbf{Question:} Which type of RAG is used in this project and why?

\begin{tcolorbox}[colback=blue!5!white,colframe=blue!75!black,title=Agentic Local RAG]
The project implements a \textbf{Stateful Agentic RAG} system using the following stack:
\begin{itemize}
    \item \textbf{Vector Store:} FAISS (Facebook AI Similarity Search).
    \item \textbf{Embeddings:} HuggingFace \texttt{all-MiniLM-L6-v2} (Local, zero cost).
    \item \textbf{Orchestration:} LangGraph for multi-agent workflows.
\end{itemize}
\end{tcolorbox}

\subsection*{The RAG Stack Summary}
\begin{table}[h]
\centering
\begin{tabularx}{\textwidth}{@{}l l X@{}}
\toprule
\textbf{Component} & \textbf{Technology} & \textbf{Reason} \\ \midrule
Loader & PyPDF / WebBaseLoader & Flexible input (PDFs, URLs, Text). \\
Splitter & RecursiveCharacterText & Maintains context by overlapping chunks. \\
Embeddings & HuggingFace (MiniLM) & Local, fast, and free. \\
Vector DB & FAISS & Industry-standard for local similarity search. \\
Orchestrator & LangGraph & Keeps knowledge (RAG store) in memory across agents. \\ \bottomrule
\end{tabularx}
\caption{Detailed RAG Stack used in the project}
\end{table}

\subsection*{Why this particular type?}
\begin{enumerate}
    \item \textbf{Zero Cost:} Local embedding and search (no per-token API fees).
    \item \textbf{Fact-Grounded:} Multi-stage retrieval reduces hallucinations.
    \item \textbf{Privacy:} Sensitive documents are processed locally.
    \item \textbf{Self-Correction:} A \texttt{refine\_node} automatically edits content for slide length.
\end{enumerate}

\newpage

\section{Topic 2: Advanced RAG Alternatives}
\textbf{Question:} What other types of RAG can be used for better accuracy?

\subsection*{Hybrid Search vs. Base Vector RAG}
\begin{table}[h]
\centering
\begin{tabularx}{\textwidth}{@{}l X X@{}}
\toprule
\textbf{Feature} & \textbf{Base Vector RAG (Current)} & \textbf{Hybrid Search} \\ \midrule
Best For & Concepts and meaning. & Precise names, acronyms, and terms. \\
Accuracy & High for general topics. & Superior for technical or niche data. \\
Complexity & Low. & Medium (Requires keyword index). \\ \bottomrule
\end{tabularx}
\caption{Comparison: Base Vector Search vs. Hybrid Search}
\end{table}

\subsection*{Comparison of Advanced RAG Architectures}
\begin{table}[h]
\centering
\begin{tabularx}{\textwidth}{@{}l l l X@{}}
\toprule
\textbf{Method} & \textbf{Effort} & \textbf{Accuracy} & \textbf{Impact on Speed} \\ \midrule
\textbf{Current (Base)} & - & Baseline & Fast \\
Hybrid Search & Low & ++ & Fast \\
Re-ranking & Medium & +++ & Slow (Extra inference step) \\
Multi-Query & Low & ++ & Medium (Multiple searches) \\
GraphRAG & High & ++++ & Slow (Graph traversal) \\ \bottomrule
\end{tabularx}
\caption{Technical Comparison of Advanced RAG Methods}
\end{table}

\newpage

\section{Topic 3: URL Data Fetching \& PPT Generation}
\subsection*{The URL Pipeline}
\begin{enumerate}
    \item \textbf{Data Fetching:} \texttt{WebBaseLoader} extracts text from URLs.
    \item \textbf{Ingestion:} Text is chunked (1,000 chars) and stored in FAISS.
    \item \textbf{Generation:}
    \begin{itemize}
        \item \textbf{Planner:} Creates an outline.
        \item \textbf{Writer:} Performs similarity search for each slide header.
        \item \textbf{Refiner:} Shortens text to fit slide layouts.
    \end{itemize}
\end{enumerate}

\subsection*{Why the Hybrid Approach (AI + Programmatic)?}
\begin{itemize}
    \item \textbf{Structural Reliability:} \texttt{python-pptx} ensures exact rendering and no text overflow.
    \item \textbf{Premium Design:} A Designer agent picks custom themes based on the topic.
    \item \textbf{Performance:} Threaded image fetching is significantly faster.
\end{itemize}

\end{document}
